\begin{table*}[!t]
\centering
\caption{Three complementary evaluation blocks. R: reported in existing formal
or diagnostic cohorts; I: implemented or exploratory, but final validation or
coverage is incomplete; P: planned. R does not mean that every model, song, and
execution layer has been evaluated. No aggregate quality score is used.}
\label{tab:evaluation_matrix}
\small
\renewcommand{\arraystretch}{1.16}
\setlength{\tabcolsep}{0pt}
\begin{tabular}{@{}>{\raggedright\arraybackslash}p{0.15\textwidth}
                    @{\hspace{0.02\textwidth}}
                    >{\raggedright\arraybackslash}p{0.36\textwidth}
                    @{\hspace{0.02\textwidth}}
                    >{\raggedright\arraybackslash}p{0.45\textwidth}@{}}
\toprule
\textbf{Block} & \textbf{Metrics and direction} & \textbf{Question and evidence status} \\
\midrule
\textbf{1. MMR} & MMR-MS, MMDist $\downarrow$; audio-to-motion and motion-to-audio R@K $\uparrow$, median/mean rank $\downarrow$ &
Does an independent learned evaluator associate the motion with its music?
\textbf{R:} FineDance matching/retrieval subsets. \textbf{I:} cross-dataset calibration;
not a direct aesthetic judgment. \\
\midrule
\textbf{2. Music adaptation} & BAS $\uparrow$; beat-event coverage and timing; onset--motion correlation and response lag &
Does the movement respond to the soundtrack's rhythm and dynamics?
\textbf{R:} BAS and beat diagnostics in formal cohorts; onset/impact diagnostics in
layer analysis. Event coverage is not a mandatory one-to-one beat assignment. \\
& Tempo/phase, section response, style/affect fit; correct versus swapped-audio preference &
\textbf{I:} tempo/style diagnostics. \textbf{P:} validated phrase/affect endpoints
and blinded cross-paired music-fit study at generation and execution. \\
\midrule
\textbf{3. Dance quality and aesthetics} & FID$_k$/FID$_g$/FID$_L$ $\downarrow$;
Div$_k$/Div$_g$/Div$_L$ and same-song multimodality $\leftrightarrow$GT &
Is the motion distribution plausible and varied? \textbf{R:} kinetic/geometric
and learned-metric subsets. Small GT cohorts limit FID interpretation. \\
& Velocity, acceleration, jerk, activity and repetition $\leftrightarrow$GT;
boundary jumps and physical violations $\downarrow$ &
\textbf{R:} motion-dynamics and continuity diagnostics. \textbf{I:} contact/skating
proxies and standardized physical checks. Low activity or low jerk alone is not success. \\
& Blinded naturalness, smoothness, expressiveness and aesthetics $\uparrow$ &
\textbf{P:} silent-video quality study and paired reference/execution study;
no perceptual superiority is claimed from existing proxy scores. \\
\bottomrule
\multicolumn{3}{@{}l@{}}{\footnotesize $\uparrow$: higher is better; $\downarrow$:
lower is better; $\leftrightarrow$GT: match the paired ground-truth distribution.} \\
\end{tabular}
\end{table*}

\subsection{Shared Evaluation Contract}
Table~\ref{tab:evaluation_matrix} separates three questions: learned cross-modal
matching (MMR), explicit music adaptation, and dance quality/aesthetics. They are
complementary, not independent quantities that can be added into one score.
The third block concerns the motion itself; the first two concern its relation
to the soundtrack. Distributional, rhythmic, physical, and perceptual evidence
follows established dance evaluation practice~\cite{li2021aist,tseng2023edge,li2024lodge,ji2026discoforcing,li2026roboperform}.

The calibration chain comprises paired human motion ($O_{\mathrm{Human}}$), its
GMR-retargeted G1 counterpart ($O_{\mathrm{G1}}$), and execution of that G1 target
($O_{\mathrm{Exec}}$). Generated motion is evaluated through the recorded
$M_{\mathrm{gen}}$--$M_{\mathrm{sent}}$--$M_{\mathrm{target}}$--$M_{\mathrm{exec}}$
chain. The generator is trained directly in the G1 domain; online output is not
retargeted again. The final protocol applies all three blocks to both
$M_{\mathrm{gen}}$ and $M_{\mathrm{exec}}$ where the evaluator's domain calibration
permits it. Existing tracking recordings do not imply that this complete scoring
matrix has already been computed.

Audio time, crop duration, rig, feature normalization, evaluator checkpoint, and
aggregation rules are fixed within a comparison. Music scores use the recorded
audio--action clock, without independently shifting each video to improve fit.
Lag-compensated tracking is a separate diagnostic. Startup is reported untrimmed
and separately from the post-startup interval. Missing inputs, unsupported
domains, and unvalidated endpoints are marked N/A, not treated as zero-quality
samples or silently removed. Music-conditioning features are not evaluation
metrics; the learned evaluator is frozen independently of generator fitting.

The source song, rather than the frame, is the statistical unit. Sampling seeds
are repeated measurements within a song, distinct from training seeds and
tracker resets. Final comparisons require a held-out manifest, all eligible
songs, paired song-level effects with 95\% confidence intervals, and tempo/style
strata based on verified metadata. Demonstration selection is separate from
quantitative inference. Different human/G1 feature spaces or evaluator versions
must not share a numerical ranking. Generated-to-GT joint RMSE is not an overall
choreography score because valid dance is one-to-many.

\subsection{Block 1: MMR Matching and Retrieval}
\textbf{Purpose and evaluator.} Music--motion retrieval (MMR) tests whether
corresponding audio and motion are closer than controlled mismatches in an
independently learned space. Our implemented evaluator uses 35-D Librosa
temporal audio descriptors and frozen MuQ-large global audio features with
native G1 motion features; its two towers
are fitted on paired training GT, then frozen for scoring. Global/local
contrastive and time-shift controls calibrate correspondence. Music identity,
not dancer or performance identity, defines positives. The model is an
evaluation instrument, not a new generator contribution or a substitute for
human judgment. Offline evaluation may inspect the complete scored segment;
these evaluator features are not supplied to the causal generator. This does
not grant future-waveform access to any deployable conditioning route.

\textbf{Paired matching.} Let $u_g,v_g$ be global music/motion embeddings and
$u_t,v_t$ their aligned, non-overlapping one-second segment embeddings. We use
the global--temporal matching distance adopted from SoulNet~\cite{li2025soulnet}:
\begin{align}
 D_{\mathrm{temp}} &= \sum_{t=2}^{T}
 \left\|(u_t-u_{t-1})-(v_t-v_{t-1})\right\|_2,\\
 d_{\mathrm{MMR}} &=
 \sqrt{0.7\|u_g-v_g\|_2^2+0.3D_{\mathrm{temp}}}.
 \label{eq:mmr_score}
\end{align}
MMR-MS averages this per-pair distance; lower is better. We also retain the
global and temporal terms to expose their different behavior. The temporal
term is a sum of norms, not squared norms or a time average; crop length must
therefore be matched. A global-only fallback is not the full MMR-MS.
MMDist separately measures mean paired global Euclidean distance; where
multiple matching candidates exist, the implementation uses their centroid.
It is not maximum mean discrepancy (MMD) and not a retrieval success rate.

\textbf{Retrieval.} We report audio-to-motion and motion-to-audio R@K for
$K\in\{1,2,3,5,10\}$, together with median and mean rank. A query succeeds if a
valid same-music candidate appears among its first $K$ results. Candidate pools,
query counts, distance convention, and multiple-positive rules remain fixed;
the first valid positive defines the rank. Chance comparisons must reflect
the actual pool and number of positives. Correct versus wrong-track and
wrong-time margins test whether the evaluator detects more than generic
movement, genre similarity, or performance identity.

\textbf{Status and limits.} FineDance-only matching/retrieval results are
available for formal generator subsets. A multi-dataset candidate has not
closed all AIST++ calibration gates, as detailed below. No cross-dataset or
executed-motion aesthetic ranking follows automatically from a lower MMR-MS.
Evaluator identity, training/test music separation, feature versions, and
temporal-embedding availability accompany each reported score.

\subsection{Block 2: Music Adaptation}
\textbf{Rhythmic alignment.} For motion-accent times $\mathcal{B}_x$ and
audio-beat times $\mathcal{B}_a$, motion-to-music Beat Alignment Score is
\begin{equation}
 \mathrm{BAS}=\frac{1}{|\mathcal{B}_x|}
 \sum_{b\in\mathcal{B}_x}
 \exp\!\left(-\frac{\min_{a\in\mathcal{B}_a}(b-a)^2}{2\sigma^2}\right).
 \label{eq:bas}
\end{equation}
The event extractor, smoothing, and $\sigma$ are frozen and recorded with the
metric version. The layer analyzer derives motion accents from minima of
smoothed mean joint speed and uses $\sigma=0.10$~s. Reverse-direction BAS is
labelled separately; historical scores require their own extractor/version
audit before pooling.
We retain event counts, beat density, nearest-event timing and offbeat
diagnostics: a high BAS with very few events does not establish sustained
rhythmic dance. Historical Beat-P/R/F1 values describe event coverage under a
specified matching tolerance; without an explicit beat-assignment protocol
they are not mandatory quality endpoints. Dance need not hit every audio beat,
and half-time or double-time phrasing is not inherently incorrect.

\textbf{Dynamics, tempo, and structure.} We measure zero-lag audio-onset
correlation with joint-speed and motion-impact curves. The implemented impact
proxy is the smoothed absolute derivative of the speed curve, not a measured
contact force. Best-lag correlation and its lag diagnose delayed response;
a lag from weak correlation is unreliable and is not used to retime the main
music score. Tempo/period and accent-phase diagnostics complement beat proximity
but require declared half/double-tempo handling. Phrase/section transition
response, annotated style fit, and affect agreement are further endpoints,
not inferred from BAS or motion energy. Validated phrase and affect evaluation
remains pending; human facial-expression metrics are not directly applicable
to G1's unactuated face.

\textbf{Appropriate adaptation rather than mere change.} Current formal
correct/null, wrong-track, wrong-time, and shifted-condition controls test
music dependence. The final cross-paired protocol additionally generates
$D_A,D_B$ from songs $A,B$ with matched motion history and sampling noise, and
scores both dances with both soundtracks. For a higher-is-better fit score $s$,
\begin{align}
 \Delta_{\mathrm{fit}}=\tfrac12\big[&s(A,D_A)-s(B,D_A)\nonumber\\
                              +&s(B,D_B)-s(A,D_B)\big].
 \label{eq:cross_music_fit}
\end{align}
For a distance metric, $s=-d$. Positive matched preference is sought together
with dance quality, not merely a large $\|D_A-D_B\|$. Same-style/similar-tempo
negatives test fine musical specificity; cross-style pairs test broader
response. This design is repeated for generated and executed dances without
changing their audio alignment. Blinded listeners rate rhythm, accent/dynamic
response, and overall music--dance fit; these perceptual controls remain TODO.

\subsection{Block 3: Dance Quality and Aesthetics}
\textbf{Distribution and diversity.} FID$_k$ and FID$_g$ compare generated and
held-out GT feature distributions using kinetic and geometric G1 descriptors,
respectively. FID$_L$ applies the same mean/covariance distance in an independent
frozen motion-embedding space. Div$_k$, Div$_g$, and Div$_L$ measure pairwise
feature distances and are interpreted relative to matched GT, not maximized
without bound. Same-song multimodality is the mean pairwise motion-embedding
distance across sampling seeds, macro-averaged over eligible songs; the formal
protocol uses at least three seeds. Here multimodality (MM) is distinct from
cross-modal MMR. Global diversity cannot show that one song admits diverse
valid dances. Small GT samples, including the current 18-song cohort, make FID
auxiliary evidence rather than a robust standalone superiority claim.

\textbf{Local and long-horizon motion.} We retain joint/root velocity,
acceleration and jerk distributions, high-frequency jitter, and position and
velocity discontinuities at C4 boundaries. Activity is measured using joint
excursion and squared angular velocity; static-window fraction and repeated
motion expose freezing or looping. These values are compared with GT and
controlled corruptions: standing still minimizes jerk but is not a good dance,
and high activity can be jitter. Reconstruction NRMSE, branch-swap sensitivity,
and continuous-code effective rank test the D+C mechanism, while the three
evaluation blocks test generated outcomes. They are not interchangeable.

\textbf{Physical plausibility.} Contact-weighted skating, penetration,
joint/dynamic-limit violations, root drift, and stability complement motion
statistics. The current G1 analyzer includes a physical-foot-contact (PFC)
proxy combining planar root acceleration and foot motion, and foot-skating
rate (FSR) proxies based on estimated contact. Their ground height, contact
threshold, speed threshold, and world/body frame must be disclosed; they are
not silently equated with standardized human-motion benchmarks. Kinematic
plausibility does not prove dynamic feasibility. Actual survival, falls, and
interventions are measured in controller execution.

\textbf{Perceptual quality.} A planned blinded silent-video study rates
naturalness, continuity/smoothness, expressiveness, and overall aesthetic
quality separately from the soundtrack-fit study. A second paired study asks
whether execution retains the reference's expressive movement. Presentation
matches rig, camera, duration, fps, and real playback speed, with randomized
order and participant/song-aware analysis. Low FID, low jerk, and high tracking
retention do not individually prove that dance is beautiful. No completed
human-quality result is claimed in this draft.

\subsection{Execution and Timing as Supporting Diagnostics}
Tracking and runtime measurements accompany, rather than replace, the three
quality blocks. We report unaligned and lag-aligned joint RMSE, root drift,
target-to-execution lag, joint-wise amplitude and velocity-power retention,
and limb/frequency-band attenuation. Startup and post-startup results are
separate. Audio cutoff, forecast publication/fetch, action deadlines, condition
age, valid FMS count, planning-time tails, deadline misses, holds/aborts, and
simulation real-time factor establish whether the intended online contract
was realized. Recording coverage and resource logs qualify comparisons;
successful retries alone cannot estimate reliability. A stable robot with
attenuated accents may track acceptably while scoring poorly on music fit.

\papertodo{T04/T08--T10: complete the three-block score matrix for generated and
executed motion; freeze event tolerances, windows, evaluator hashes, and
small-sample reporting. Run the cross-paired music-fit test and blinded aesthetic
study. Exploratory tempo/style/contact metrics and planned phrase/affect metrics
must not be presented as fully validated results.}

